\section{Terminal theorem and corollaries}
\label{sec:terminal}

\begin{theorem}[Exhaustion of same-domain incompleteness objection routes]
\label{thm:incompleteness-routes}
Every same-domain incompleteness objection route falls into one of the following six families:
\begin{enumerate}[label=(\roman*), leftmargin=2.6em]
    \item code-level readback routes;
    \item semantic or intended-model rescue routes;
    \item reflection, infinitary, or transfinite closure routes;
    \item second-equivalence or hidden-refinement routes;
    \item richer same-domain extension or new-gate routes; and
    \item explicit scope-change routes.
\end{enumerate}
The first five are impossible or closure-inert inside the coherent-reasoning kernel; the sixth is not same-domain.
\end{theorem}

\begin{proof}
Code-level readback routes are exhausted by \cref{thm:code-role-exhaustion}. Any route that genuinely threatens closure must also satisfy the package in \cref{thm:godel-prerequisites}; \cref{thm:no-diagonalization} then closes the negative fixed-point branch. Semantic, reflection, infinitary, and meta-foundational rescue routes are exhausted by \cref{thm:rescue-exhaustion}. Any surviving second-equivalence or richer same-domain extension route would reintroduce hidden gate work or a new same-domain authority source, again excluded by \cref{thm:code-role-exhaustion,thm:rescue-exhaustion}. What remains is explicit scope change, which is not a same-domain threat by definition.
\end{proof}

\begin{theorem}[Terminal non-instantiability theorem]
\label{thm:main}
No G\"odel-style same-domain threat architecture is instantiable in the unique admissible interior of the fixed coherent-reasoning regime.
\end{theorem}

\begin{proof}
Suppose such an architecture were instantiable. Then it would belong to one of the six route families of \cref{thm:incompleteness-routes}. The first five are impossible or closure-inert in the fixed kernel. The sixth is explicit scope change and therefore not same-domain. This contradiction shows that no same-domain G\"odel threat architecture is instantiable.
\end{proof}

\begin{corollary}[No same-domain ``true but unprovable'' split with governance force]
There is no same-domain split between a standing-bearing content and its internal provability-like classification that survives as a governance-relevant closure threat on the fixed domain.
\end{corollary}

\begin{proof}
Any such split would require a same-domain G\"odel threat architecture and is excluded by \cref{thm:main}.
\end{proof}

\begin{corollary}[Kernel-relative reading]
Theorem~\ref{thm:main} is a kernel-relative non-instantiability result. It leaves classical incompleteness untouched while eliminating same-domain threat architectures inside the fixed coherent-reasoning regime.
\end{corollary}
